\documentclass[12pt]{article}
\usepackage[utf8]{inputenc}
\usepackage[T1]{fontenc}
\usepackage{hyperref}
\title{Ixtayutla Mixtec Language Resources for K-12 Curriculum Development}
\author{Compiled by Claude Sonnet 5 (claude-sonnet-5) for Andr\'as Kornai}
\date{}
\begin{document}
\maketitle

\section{Introduction}
Ixtayutla Mixtec (ISO 639-3 vmj; Glottocode ixta1235) is an Otomanguean language of Oaxaca, Mexico. Searches found very little dedicated material. No school textbook, primer, grammar, or dictionary specific to this variety turned up. The resources below are general aggregator, scripture, and audio pages only.

\section{Primary Educational Resources}
No primers, textbooks, or school programmes for Ixtayutla Mixtec were found in search results.

\section{Linguistic and Reference Materials}
\begin{itemize}
\item Wikipedia stub on Ixtayutla Mixtec: \url{https://en.wikipedia.org/wiki/Ixtayutla_Mixtec}
\item Glosbe vmj--English word and phrase pages. These come from parallel corpora and community contributions and are not authoritative: \url{https://glosbe.com/vmj/en}
\end{itemize}

\section{Religious and Community Resources}
\begin{itemize}
\item ScriptureEarth page with Bible text and audio: \url{https://www.scriptureearth.org/00eng.php?iso=vmj}
\end{itemize}

\section{Digital and Audio Resources}
\begin{itemize}
\item Global Recordings Network audio in Ixtayutla Mixtec (evangelistic and Bible-teaching content): \url{https://www.globalrecordings.net/en/language/vmj}
\end{itemize}

\section{Curriculum Development Context}
Little is documented online. A curriculum would have to rely on community work and original documentation. Academic sources such as SIL publications or the Glottolog bibliography may hold more, but they were not verified here.
\end{document}